\documentclass[../cuaderno6.tex]{subfiles}

\begin{document}

\subsection{Cubo}
En el caso del cubo, el campo coercitivo $H_c$, la magnetización remanente $M_r$ y la magnetización máxima $M_{max}$ no varian con el ángulo. Se puede ver en las figuras \ref{fig:H_c_segun_H_a_Cubo_Fa0_29_09} y \ref{fig:H0_c_angulo_minors_Cubo} como para cada ángulo medido, las mediciones casi no difieren entre sí. El ajuste para el campo coercitivo tampoco difiere de forma significativa. Rotar esta pieza no genera diferencias significativas en las mediciones del campo coercitivo, magnetización remanente ni máxima. Tampoco se observan diferencias mínimas dadas por la estructura interna del material, que podría ser isótropa.

\begin{figure}[h]
    \centering
    \includegraphics[width=0.5\textwidth]{../img/Fa0/Cubo/H_c_segun_H_a_Cubo_Fa0_29_09.png}
    \caption{El campo coercitivo $H_c$ en función del logaritmo del campo aplicado $H_{max}$ para cada ángulo.}
    \label{fig:H_c_segun_H_a_Cubo_Fa0_29_09}
\end{figure}
\begin{figure}[h]
    \centering
    \includegraphics[width=0.5\textwidth]{../img/Fa0/Cubo/H0_c_angulo_minors_Cubo.png}
    \caption{Gráfico de la pendiente del ajuste $H_c^0$ en función del ángulo.}
    \label{fig:H0_c_angulo_minors_Cubo}
\end{figure}

Al calcular el factor de desmagnetización $N_d$ del cubo, observamos a partir de la figura \ref{fig:N_d_angulo_major_Cubo} que este sí mantiene una dependencia visible, pero no relacionada a la geometría de la pieza al no respetar las simetrías de la misma con su variación. Vemos que resulta posible identificar una dependencia angular en las mediciones mediante un análisis más fino, ya que en este caso las mismas presentan diferencias a partir de su tercera cifra significativa. También es por esta razón que consideramos las mediciones obtenidas para el cubo resultan invariantes frente a rotaciones en la gran mayoría de su carácter.
    \vspace{1.5em}
\begin{figure}[H]
    \centering
    \includegraphics[width=0.5\textwidth]{../img/Fa0/Cubo/N_d_angulo_major_Cubo.png}
    \caption{Factor de desmagnetización $N_d$ en función del ángulo.}
    \label{fig:N_d_angulo_major_Cubo}
\end{figure}

\end{document}